% =================================================================
% CAPÍTULO 5: DETALHES DE IMPLEMENTAÇÃO (TEMPLATE)
% Objetivo: registrar decisões técnicas e arquitetura da solução.
% Dicas:
%  - Mantenha diagramas atualizados com o texto.
%  - Use labels descritivos (ex.: fig:diagrama1, lst:prototipo)
% =================================================================
\section{Detalhes de Implementação}

[Descreva detalhes técnicos relevantes da implementação e do design.]

\subsection{Diagrama(s) de Arquitetura}
\begin{figure}[h]
  \centering
  \includegraphics[width=\linewidth]{template_diagram1.png}
  \caption{[Legenda do diagrama 1]}
  \label{fig:diagrama1}
\end{figure}

\begin{figure}[h]
  \centering
  \includegraphics[width=\linewidth]{template_diagram2.png}
  \caption{[Legenda do diagrama 2]}
  \label{fig:diagrama2}
\end{figure}

\subsection{Listagens de Código (Opcional)}
% \begin{lstlisting}[style=code, caption={<Legenda do trecho de código>}, label={lst:<identificador>}]
% <Exemplo de código ou pseudocódigo>
% \end{lstlisting}
[Inclua listagens de código, se aplicável ao trabalho.]
